\chapter{Injection beam loss event}\label{sec:Injection beam loss event}

To further increase the luminosity towards its target value, the currents of both the electron and positron need to be increased by approximately factor of 2. One of the obstacles to increase the beam currents is injection beam loss. During injection, injected bunches can deviate from their ideal trajectories and strike collimators, the beam pipe, or detector components, producing radiation showers. These showers can degrade data acquisition (DAQ) performance and expose the detectors to radiation leveles approaching their tolerance limits. Beam losses at the collimators also hinder further increases in beam current. 